\section{Introduction}

Aquatic ecosystems face pressure from chemical pollution. There are hundreds of thousands of chemicals in use in society \cite{wang2020toward}, and many of them find their way into surface water systems. There are increasing regulatory efforts to prevent and control this, ranging from the recently proposed \textit{Safe \& Sustainable by Design} \cite{caldeira2022safe} framework as a pro-active approach in the design of novel chemicals, up to the retrospective water quality assessment and mitigation of the EU Water Framework Directive (2000/60/EC). Reaching the societal goal of \textit{sufficient water of good quality}, for the many human uses of clean water and as habitat for aquatic species assemblages, requires the risk assessment and management of individual chemicals and chemical mixtures.

One of the key requirements in risk assessment and management is to know the hazards of the chemicals that we use. That knowledge is commonly collected by performing single-species, single-compound ecotoxicity tests in a laboratory setting. For a given chemical, standardised tests are typically performed on a handful of selected species, such as an algae, a daphnid and/or a fish. The purpose of these tests is to determine the concentration at which exposure is lethal (LCx), or has a specific negative effect on a life history characteristic (ECx), such as growth or reproduction, for x\% of the test individuals.

The results and conditions of these tests are being collated and curated by several organisations. For example, the US EPA ECOTOX Knowledgebase now contains over a million test records spanning thousands of species and chemicals \cite{ECOTOX}. However, this still represents a mere fraction all species--chemical pairs.

Due to this data sparsity, we typically only know at which concentration level a chemical is toxic for a minority of species. The traditional approach to deal with this is to fit a parametric (typically lognormal) distribution on the available single-species toxicity values of a given chemical, the so-called Species Sensitivity Distribution (SSD). This describes the sensitivity of an entire ecological community to a chemical \cite{posthuma2001species}, and can be used to derive regulatory thresholds such as the Hazardous Concentration for 5\% of species (HC5). However, this fit is only reliable when the number and variety of tested species is sufficient, which is only the case for a small fraction of all chemicals.

The growing data availability, together with the sparsity problem, has motivated a wave of machine learning approaches that use the collated testing data on known species--chemical pairs to predict the toxicity of untested chemical--species pairs, or even of entirely new chemicals or never-tested species. \citet{von2023potential} identify key parameters where machine learning could address data gaps in toxicity characterisation\iffalse OL: Can we be more specific here? What were those key parameters?\fi, \citet{kosnik2024harnessing}, advocate for computational methods as essential to characterising chemical impacts on biodiversity at scale, while \citet{dong2025machine} review the expanding role of machine learning across exposure estimation, biotoxicity prediction, and environmental fate modelling.

A more critical account is provided by \citet{schur2025comparability}, who highlight significant comparability issues in predictive ecotoxicology, calling for standardised benchmarks and evaluation protocols. To address this, the authors have curated the ADORE collection of benchmark datasets, based on the US EPA ECOTOX Knowledgebase and enriched with chemical and taxonomical features \cite{ADORE}.

The ADORE collection has enabled the systematic evaluation of various models. \citet{gasser2024fish} found that model performance was highly dependent on the data split while benchmarking tree-based and Gaussian-process regression models on the fish-specific ADORE challenge, and applying them to fish SSD construction. A novel approach was taken by \citet{Viljanen2023}, who framed ecotoxicity prediction as a recommender-system problem and showed that factorisation machines could match more complex methods in the context of highly sparse test data. This was further elaborated by \citet{PosthumaImprovingHazard}, who scaled the approach to predict LC50 values for over four million species--chemical data gaps and used the resulting predictions to construct Species Sensitivity Distributions and Chemical Hazard Distributions for all chemicals and species, representing the current state of the art for machine-learning-based SSD construction at scale. Others have explored meta-learning to share knowledge across low-resource species \cite{schlender2023bigger}, and Bio-QSARs that incorporate physiological traits to enable genuine cross-species extrapolation \cite{zubrod2024bio}.

While the overall predictive accuracy of such models can be measured using approaches like cross-validation, yielding a single accuracy metric across the entire dataset, this tells us little about how much to trust an \textit{individual} prediction (and by extension resulting scatter-SSDs and derived protective thresholds). This is because the  existing models produce only point estimates, and do not quantify the confidence of individual predictions. This is a significant limitation for practical risk assessment, where knowing \textit{how much} to trust a prediction is as important as the prediction itself. In practice, we expect prediction quality to be impacted by two sources of uncertainty. Firstly, a chemical that has been tested on dozens of species should yield more reliable predictions than one with only a handful of observations, and similarly for well-studied versus rarely tested species. Secondly, the currently available data clearly show, for species--chemical pairs that have been tested multiple times, that repeated testing of the chemical on the same species yields highly variable results. These two sources of uncertainty are commonly referred to as epistemic (arising from limited data constraining the model) and aleatoric (inherent to the biological and experimental variability of the testing process) \cite{epistemicaleatoric}.

Bayesian modelling offers a natural framework to deal with uncertainty: by treating model parameters as distributions rather than fixed values, it produces not just a single prediction but a range of plausible outcomes, directly expressing the model's uncertainty. Bayesian approaches have recently been applied to uncertainty-aware human toxicity prediction \cite{von2026uncertainty} and to hierarchical SSD construction for microplastic particles \cite{iwasaki2026bayesian}, but have not yet been applied to the large-scale species--chemical matrix completion problem in ecotoxicology.

The present work addresses this gap by extending the factorisation machine architecture applied to ecotoxicity prediciton in \citet{PosthumaImprovingHazard} into a fully Bayesian framework, the Bayesian Factorisation Machine (BFM), adapted from \citet{rendle2011bayesian}. In addition, we explicitly model how noisy the underlying data is for each chemical, allowing us to distinguish between aleatoric and epistemic uncertainty.

Thus, the BFM model presented in this work allows risk assessors to see predicted toxicity values for all combinations of species and chemicals that are present in the training data, the uncertainty surrounding these predictions, as well as to how much of this uncertainty is inherent experimental noise, and how much is due to a lack of data.

Furthermore, we also describe how this uncertainty propagates to the next steps of practical assessments. For each chemical we derive an SSD that is made up of the point-predictions for each species, along with an uncertainty band showing the space of plausible SSD curves given the model uncertainty, allowing for more informed risk assessment and regulation.

In this study, the LC50 endpoint is used as a methodological prototype: acute lethality is a relatively well-defined test endpoint (compared to e.g.\ chronic no-observed effect concentrations for life-history characteristics such as growth and reproduction) and is numerically abundant in the data, making it ideal for developing and validating the uncertainty quantification approach. However, with additional tuning and validation, the method can easily be adapted to other test endpoints such as NOEC or EC10 under chronic exposure conditions, which are often preferred as basis for regulatory applications (such as setting protective concentration standards).